\documentclass[11pt,a4paper]{article}
\usepackage[utf8]{inputenc}
\usepackage[T1]{fontenc}
\usepackage[french]{babel}
\frenchbsetup{StandardLists=true}
\usepackage[left=2cm,right=2cm,top=2cm,bottom=2cm]{geometry}
\usepackage[dvipsnames]{xcolor}
\usepackage{fouriernc}
\usepackage{amsmath,amssymb,amsfonts}
\usepackage{array,tabularx,booktabs,multirow}
\usepackage{colortbl}
\usepackage{graphicx}
\usepackage{mdframed}
\usepackage{enumitem}
\usepackage{fancyhdr}
\usepackage{chemformula}
\usepackage{awesomebox}
\setlength{\aweboxvskip}{0mm}
\usepackage{tcolorbox}
\tcbuselibrary{skins}
\usepackage{fontawesome5}

\definecolor{exoheader}{HTML}{1E5AA0}
\definecolor{exolight}{HTML}{DCE8FF}
\definecolor{warnorange}{RGB}{220,100,0}
\definecolor{problemebg}{HTML}{FFFDE7}

\renewcommand{\arraystretch}{1.8}
\setlength{\extrarowheight}{1mm}

\pagestyle{fancy}
\renewcommand\headrulewidth{1pt}
\fancyhead[L]{Académie de Besançon}
\fancyhead[C]{\textbf{TP --- Combustion d'un hydrocarbure}}
\fancyhead[R]{Terminale / BTS Chimie}
\fancyfoot[C]{page \thepage}
\fancyfoot[R]{2025 -- 2026}
\fancyfoot[L]{Alexis RUIZ}

\newcommand{\sectitle}[1]{%
  \vspace{6pt}
  \noindent\colorbox{exoheader}{\parbox{\dimexpr\linewidth-2\fboxsep}%
    {\color{white}\bfseries\large\strut\quad #1}}
  \vspace{4pt}
}
\newcommand{\subsectitle}[1]{%
  \vspace{4pt}
  \noindent\colorbox{exolight}{\parbox{\dimexpr\linewidth-2\fboxsep}%
    {\color{exoheader}\bfseries\strut\quad #1}}
  \vspace{2pt}
}
\newcounter{question}
\newcommand{\question}{%
  \stepcounter{question}%
  \noindent\textcolor{exoheader}{\textbf{Question \thequestion.}}\quad
}

\begin{document}

\noindent\textbf{NOM :}\hfill\textbf{Prénom :}\hfill\textbf{Classe :}\hfill\phantom{x}\\[0.2cm]
\hrule\vspace{0.2cm}

\begin{center}
  {\large\bfseries\textcolor{exoheader}{Travaux Pratiques de Chimie}}\\[4pt]
  {\LARGE\bfseries Combustion d'un hydrocarbure --- Pouvoir calorifique}\\[2pt]
  {\small Académie de Besançon --- Belfort}
\end{center}
\hrule\vspace{6pt}

\begin{tcolorbox}[enhanced, colback=problemebg, colframe=exoheader, boxrule=1.5pt,
  title={\faFlask\quad Problématique},
  fonttitle=\bfseries\color{white},
  attach boxed title to top left={yshift=-2mm, xshift=4mm},
  boxed title style={colback=exoheader, colframe=exoheader}]
\textit{La combustion des hydrocarbures constitue la principale source d'énergie dans les transports.
\textbf{Quelle quantité d'énergie libère la combustion d'un hydrocarbure ?
Comment la mesurer expérimentalement ?}}
\end{tcolorbox}

\vspace{6pt}

% ══════════════════════════════════════════════════════════════════════════════
\sectitle{Partie I --- Étude théorique}
% ══════════════════════════════════════════════════════════════════════════════

\subsectitle{1. Les hydrocarbures}

\begin{mdframed}[backgroundcolor=exolight, linecolor=exoheader, linewidth=1.2pt]
Les \textbf{hydrocarbures} sont des composés organiques ne contenant que du carbone et de l'hydrogène.
La \textbf{combustion complète} d'un hydrocarbure produit uniquement \ch{CO2} et \ch{H2O}.
Le \textbf{pouvoir calorifique inférieur (PCI)} est l'énergie libérée par la combustion
d'une mole (ou d'1~kg) de combustible (eau produite à l'état vapeur).
\end{mdframed}

\vspace{4pt}
\question Écrire l'équation de combustion complète du propane (\ch{C3H8}). \dotfill

\vspace{0.3cm}\noindent\dotfill

\question Écrire l'équation de combustion complète du butane (\ch{C4H10}). \dotfill

\vspace{0.3cm}\noindent\dotfill

\question Pour une bougie (paraffine, assimilée au pentacosane \ch{C25H52}),
écrire l'équation de combustion complète. \dotfill

\vspace{0.3cm}\noindent\dotfill

\question Qu'est-ce qu'une combustion \textbf{incomplète} ? Donner les produits formés.
Quels sont les dangers ? \dotfill

\vspace{0.3cm}\noindent\dotfill

% ══════════════════════════════════════════════════════════════════════════════
\sectitle{Partie II --- Protocole expérimental}
% ══════════════════════════════════════════════════════════════════════════════

\begin{mdframed}[backgroundcolor=orange!15, linecolor=warnorange, linewidth=2pt]
\textbf{\textcolor{warnorange}{\faExclamationTriangle\quad ATTENTION}} ---
Manipulation avec une flamme. Travailler loin des produits inflammables.
Porter \textbf{lunettes de protection}. Ne jamais laisser la flamme sans surveillance.
\end{mdframed}

\vspace{4pt}
\subsectitle{1. Mesure du pouvoir calorifique d'une bougie}

\textbf{Principe :} La chaleur dégagée par la combustion de la bougie est absorbée
par une masse $m_{\text{eau}}$ d'eau. On mesure l'élévation de température $\Delta T$.
\[
  E = m_{\text{eau}} \times c_{\text{eau}} \times \Delta T
  \quad \text{avec } c_{\text{eau}} = 4{,}18~\text{J}\cdot\text{g}^{-1}\cdot\text{°C}^{-1}
\]

\vspace{4pt}
\begin{enumerate}[leftmargin=1.5cm, label=\textcolor{exoheader}{\textbf{\arabic*.}}]
  \item Peser la bougie avant la manipulation : $m_1 = $ \dotfill~g
  \item Mesurer $m_{\text{eau}} = 200$~g d'eau dans le bécher en inox.
  \item Relever la température initiale de l'eau : $T_1 = $ \dotfill~°C
  \item Allumer la bougie et placer le bécher au-dessus (distance de 3~cm environ).
  \item Agiter doucement et mesurer la température toutes les minutes pendant 10~min.
  \item Éteindre la bougie. Peser immédiatement : $m_2 = $ \dotfill~g
  \item Relever la température finale : $T_2 = $ \dotfill~°C
\end{enumerate}

\vspace{4pt}
\begin{center}
\begin{tabularx}{\linewidth}{|>{\centering\arraybackslash\columncolor{exolight}}m{2cm}
                              |>{\centering\arraybackslash}m{2cm}
                              |>{\centering\arraybackslash\columncolor{exolight}}m{2cm}
                              |>{\centering\arraybackslash}m{2cm}
                              |>{\centering\arraybackslash\columncolor{exolight}}m{2cm}
                              |>{\centering\arraybackslash}X|}
\hline
\rowcolor{exoheader}
\color{white}$t$ (min) & \color{white}$T$ (°C) &
\color{white}$t$ (min) & \color{white}$T$ (°C) &
\color{white}$t$ (min) & \color{white}$T$ (°C) \\
\hline
0 & & 4 & & 8 & \\
\hline
1 & & 5 & & 9 & \\
\hline
2 & & 6 & & 10 & \\
\hline
3 & & 7 & & & \\
\hline
\end{tabularx}
\end{center}

% ══════════════════════════════════════════════════════════════════════════════
\newpage
\sectitle{Partie III --- Exploitation des résultats}
% ══════════════════════════════════════════════════════════════════════════════

\subsectitle{1. Calculs}

\question Calculer la masse de bougie consommée :
\[
  \Delta m = m_1 - m_2 = \dotfill \text{ g}
\]

\question Calculer l'énergie $E$ transférée à l'eau :
\[
  E = m_{\text{eau}} \times c_{\text{eau}} \times (T_2 - T_1) = \dotfill \text{ J}
\]

\question Calculer le PCI expérimental de la paraffine (en J/g puis en kJ/mol) :
\[
  \text{PCI} = \frac{E}{\Delta m} = \dotfill \text{ J/g} = \dotfill \text{ kJ/mol}
\]

\question Comparer avec la valeur théorique (PCI paraffine $\approx$ 46~kJ/g).
Expliquer les sources d'erreur. \dotfill

\vspace{0.3cm}\noindent\dotfill

\subsectitle{2. Identification des produits de combustion}

\question Décrire le test permettant de mettre en évidence \ch{CO2} et \ch{H2O}
parmi les produits de combustion. Indiquer les résultats attendus.

\vspace{4pt}
\begin{center}
\begin{tabularx}{\linewidth}{|>{\bfseries\columncolor{exolight}}m{3cm}
                              |>{\centering\arraybackslash}m{4cm}
                              |>{\centering\arraybackslash}X|}
\hline
\rowcolor{exoheader}
\color{white}\textbf{Produit} & \color{white}\textbf{Réactif test} & \color{white}\textbf{Résultat} \\
\hline
\ch{CO2} & & \\[0.6cm]
\hline
\ch{H2O} & & \\[0.6cm]
\hline
\end{tabularx}
\end{center}

% ══════════════════════════════════════════════════════════════════════════════
\sectitle{Annexe --- Données}
% ══════════════════════════════════════════════════════════════════════════════

\begin{mdframed}[backgroundcolor=exolight, linecolor=exoheader, linewidth=1.2pt]
\begin{itemize}
  \item Capacité thermique massique de l'eau : $c_{\text{eau}} = 4{,}18$~J$\cdot$g$^{-1}\cdot$°C$^{-1}$
  \item Paraffine (assimilée au pentacosane \ch{C25H52}) : $M = 352{,}7$~g/mol ; PCI $\approx$ 46~kJ/g
  \item Propane \ch{C3H8} : $M = 44$~g/mol ; PCI $= 46{,}4$~kJ/g
  \item Butane \ch{C4H10} : $M = 58$~g/mol ; PCI $= 45{,}7$~kJ/g
  \item Masses molaires : C = 12 ; H = 1 ; O = 16 (g/mol)
\end{itemize}
\end{mdframed}

\end{document}
